\documentclass[a4paper, 10pt]{article}
\usepackage[unicode]{hyperref}
\usepackage[utf8x]{inputenc}
\usepackage[IL2]{fontenc}
\usepackage[slovak]{babel}
\usepackage[left=2cm, top=3cm, text={17cm, 24cm}]{geometry}
\usepackage{verbatim}
\usepackage{enumitem}
\usepackage{graphicx}
\usepackage{listings} % vkladanie obrazkov
\usepackage{color}
\usepackage{xcolor}


\begin{document}
    %%%%%%%%%%%%%%%%%%%%%%%%%%%%%%%% Titulná stránka %%%%%%%%%%%%%%%%%%%%%%%%%%%%%%%%
    \begin{titlepage}
        \begin{center}
            \includegraphics[width=0.77\linewidth]{src/FIT_logo} \\
            \vspace{\stretch{0.382}}
            \Huge{Modelování a simulace -- IMS} \\
            \huge{Okruh 12. -- SHO vo výrobe} \\
            \Large{Dokumentácia}
            \vspace{\stretch{0.618}}
        \end{center}

        \begin{flushleft}
            \Large{Lapeš Zdeněk (xlapes02)} \\
            \Large{Bínovský Andrej (xbinov00)}
        \end{flushleft}
        \vspace{-12mm}
        \hfill\Large{\today}
    \end{titlepage}


    %%%%%%%%%%%%%%%%%%%%%%%%%%%%%%%% Obsah %%%%%%%%%%%%%%%%%%%%%%%%%%%%%%%%
    \tableofcontents
    \newpage

    %%%%%%%%%%%%%%%%%%%%%%%%%%%%%%%% Úvod %%%%%%%%%%%%%%%%%%%%%%%%%%%%%%%%
    \section{Úvod}
        V tejto práci je riešená implementácia procesu výroby chleba, ktorá je použitá na zostavenie modelu testujúceho
        najlepšie možné nakonfigurovanie množstva strojov, pracovníkov a miestností výroby pre rôzne scenáre požiadaviek
        množstva chlebov v daný deň. Danú prácu vypracovali študenti Zdenek Lapeš a Andrej Bínovský z Fakulty informačných
        technológií VUT v Brne.
        \subsection{Čerpanie a konzultácia dát}
            V rámci tejto práce boli dáta použité a konzultované z praxe, ktoré boli získané z výrobného procesu chleba
            z Brnenskej pekárne \texttt{Crocus}\cite{Crocus}. Pre dosianutie čo najvalidnejšieho modelu procesu výroby chleba sme
            spomenutú pekáreň navštívili a získali informácie o problematike každého procesu výroby.
        \subsection{Overovanie validity dát}
            Po celý čas modelovania boli dáta overované a validované na základe komunikácie majitela pekárne. Hned zo
            začiatku sme si namodelovali a porovnali presnú konfiguráciu spomenutej pekárne. Čo nám docielilo overenie
            presnosti modelu na základe porovnania skutočného času výroby s časom výroby chleba v modeli.


    %%%%%%%%%%%%%%%%%%%%%%%%%%%%%%%% Rozbor témy a použitých metód/technológií %%%%%%%%%%%%%%%%%%%%%%%%%%%%%%%%
    \section {Rozbor témy}
        Postup výroby chleba sa skladá z viacerych procesov, ktoré sú navzájom závislé:
        \begin{itemize}
            \item \textbf{Výroba cesta} -- Proces výroby cesta spočíva zo zmiešania všetkých surovín do jedného. Počet
            a dostupnosť surovín sa v modeli neberie k úvahe. Jeden proces výroby cesta reprezentuje jeden mixér a
            výsledokm je \textbf{140kg cesta}. Proces trvá \textbf{10 minút} \textpm \textbf{ 2 minúty}.

            \item \textbf{Krájanie cesta na bochníky} -- Krájanie cesta na bochníky je proces, ktorý sa vykonáva pracovníkom
            ručne. V modeli teda platí, že jeden pracovník sa rovná jednému stolu na krájanie. Vstupom procesu je
            \textbf{1kg cesta} a výstupom je \textbf{1 bochník}. Proces trvá \textbf{30 sekúnd} \textpm \textbf{ 5 sekúnd}.

            \item \textbf{Fermentácia bochníkov} -- Po nakrájaní sa bochníky ukladajú na plech do vozíka.
            Jeden vozík obsahuje miesto na \textbf{70 bochníkov}. Po naplnení sa vozík odvezie do fermentačnej miestnosti
            na čas  \textbf{20 minút} \textpm \textbf{ 2 minúty}.

            \item \textbf{Pečenie chleba} -- Po fermentácii sa bochníky na vozíkoch uložia do pece, kde sa pečú.
            Jedna pec má kapacitu \textbf{1 vozíku}, teda \textbf{70 bochníkov}. Proces trvá \textbf{30 minút} \textpm \textbf{ 2 minúty}.

            \item \textbf{Balenie chleba} -- Po upečení pracovníci roztrieda chleba do bedien. Balenie jedného chleba
            praconíkom trvá \textbf{10 sekúnd} \textpm \textbf{ 3 sekundy}.
        \end{itemize}

        \subsection{Použité metódy, postupy a technológie}
        Na implementáciu modelu bol použitý programovací jazyk \texttt{C++} za podpory simulačnej knižnice \texttt{SimLib}.
        Hlavným dôvodom použitia jazyka \texttt{C++} je využitie naimplementovanej
        knižnice \texttt{SimLib}. Knižnicu sme si vybrali z dôvodu modelovania procesov výroby na základe petriho siete, ktorá
        toto modelovanie podporuje. Príklady a správne použitie knižnice sme čerpali z pradnášok predmetu \texttt{IMS}\cite{IMS_slides}.
        Na kompiláciu sme využili nástroj \texttt{CMake}. Na vývoj bol zvolený operačný sýstém
        \texttt{Linux Ubuntu}, ktorý bežal vo virtualizovanom prostredí nástroja \texttt{Docker}.

        \subsection{Pôvod metód a technológií}
    \begin{itemize}
        \item \textbf{C++} -- Verzia \texttt{C++20}: \texttt{https://en.cppreference.com/w/cpp/20}.
        \item \textbf{SimLib} -- Verzia \texttt{3.09}\cite{simblib}. Autori knižnice sú Petr Peringer,
        David Leska a David Martinek.
        \item \textbf{CMAKE} -- Verzia \texttt{3.5}: \texttt{https://cmake.org/}.
        \item \textbf{Docker} -- Verzia \texttt{20.10.17}: \texttt{https://www.docker.com/}.
        \item \textbf{Linux Ubuntu} -- Verzia \texttt{20.04}: \texttt{https://ubuntu.com/}.
    \end{itemize}

    %%%%%%%%%%%%%%%%%%%%%%%%%%%%%%%% Implementácia programu %%%%%%%%%%%%%%%%%%%%%%%%%%%%%%%%
    \section {Koncepcia modelu}
    Konceptuálny model pekárne je zobrazený pomocou petriho siete\ref{fig:conceptual_model}.
    V modeli je zobrazený zjednodušený proces výroby chleba. Napriek zjednodušeniu modelu je zachovaná korektnosť simulácie.
    \begin{itemize}
        \item \textbf{Výroba cesta} -- Pri prvom procese je zanedbateľná dostupnosť a naskladnovaie surovín,
        pretože to zaobstaráva externá firma. Výstup a čas procesu je totožný s realitou.

        \item \textbf{Krájanie cesta na bochníky} -- Pri druhom procese sme reálne dáta dekomponovali. Získané
        dáta z pekárne reprezentovali čas nakrájania 70 bochníkov (zo 70 kg cesta) na jeden vozík pracovníkmi za čas \textbf{10 min} \textpm \textbf{ 2 minúty}.
        Pre systém sme zvolili proces krájania 1kg cesta na 1 bochník, 1 pracovníkom za čas \textbf{30 sekúnd} \textpm \textbf{ 5 sekúnd}.

        \item \textbf{Fermentácia bochníkov} -- Pri procese fermentácie sme zachovali vstupy, výstupy a časy na základe reálnych dát.

        \item \textbf{Pečenie chleba} -- Pri procese pečenia sme zachovali vstupy, výstupy a časy na základe reálnych dát.

        \item \textbf{Balenie chleba} -- Proces balanie chleba reprezentuje viacero podprocesov ako napríklad príprava a čistenie bedien.
        Balenie spočíva v roztriedení chleba do bedien a ich následne uloženie do skladu. Reálny proces bol meraný na
        zabalenie jedného vozíku 70 chlebov do bedien za čas \textbf{10 min} \textpm \textbf{ 2 minúty}. Jedna bedna môže obsahovať 3 až 4 chleby. Proces v modeli sme dekomponovali
        na zabalenie jedného chleba jedným pracovníkom za čas \textbf{10 sekúnd} \textpm \textbf{ 3 sekundy}.
    \end{itemize}

    \section{Architektúra simulačného modelu}

    \section{Podstata imulačných experimentov a ich priebeh}

    \section{Zhrnutie simulašnych experimentov}


    \newpage
    \bibliographystyle{czechiso}
    \bibliography{documentation}
\end{document}

